\section{La idea}\label{sec:idea}

\subsection{Para qué existe KalaScript}
KalaScript es mi escritura para tomar apuntes a mano. La inventé porque me
gusta cómo la matemática dice mucho con pocos signos: \(\forall\), \(\exists\),
\(\in\), \(\Rightarrow\). Quiero esa misma claridad en los apuntes de cualquier curso:
computación, matemática, historia, psicología, filosofía, química y lo que venga.

No intenta reemplazar el español ni la notación de cada ciencia. Es una capa
de signos que se mezcla con ellos y se ocupa de lo que los apuntes normales
dicen mal: \emph{cómo se relacionan las ideas}. En un cuaderno común todo se
une con la misma flecha «\(\to\)», y una semana después ya no sé si quise decir
«causa», «después», «implica» o «se convierte en». En KalaScript cada una de
esas relaciones tiene su propio signo, y la flecha queda solo para cuando no
importa cuál es.

\subsection{Cuatro capas en la misma línea}
Un apunte en KalaScript mezcla libremente:
\begin{enumerate}
\item \textbf{Español telegráfico} para lo que no tiene signo: nombres propios,
 términos técnicos, datos (secciones~\ref{sec:telegrafico} y~\ref{sec:abreviaturas}).
\item \textbf{Símbolos matemáticos} que ya conozco, con un significado fijo:
 \(\neg\), \(\in\), \(\subset\), \(\to\), \(\Rightarrow\), \(\therefore\), \(\uparrow\), \(\downarrow\)
 (sección~\ref{sec:matematica}).
\item \textbf{Veintinueve signos} para conectar ideas, clasificar nombres
 y marcar el margen (sección~\ref{sec:signos}).
\item \textbf{Veintinueve raíces} que se combinan para nombrar conceptos:
 mente, poder, vida, cambio\ldots{} (sección~\ref{sec:raices}).
\end{enumerate}

\begin{apunte}[Psicología · un ejemplo de las tres capas]
\linea{}{\ks{K16} encuesta 2024: horas de pantalla \ks{K28} ansiedad}
\linea{\ks{K31}}{\ks{K28} \(\neq\) \ks{K23}}
\linea{\ks{K33}}{\(\downarrow\) sueño \ks{K23} \(\uparrow\) ansiedad}
\end{apunte}
\lectura{Según una encuesta de 2024, las horas de pantalla y la ansiedad van
juntas. Importante: que vayan juntas no significa que una cause la otra.
Hipótesis: dormir menos causa más ansiedad.}

\subsection{Principios}
\begin{enumerate}
\item \textbf{Primero la mano.} Un signo existe si lo escribo rápido, lo
 distingo de los demás cuando escribo feo y lo entiendo al releer.
\item \textbf{Solo donde ayuda.} Si la palabra en español es igual de rápida y
 clara, uso la palabra. «y», «o» y «no» no necesitan signo propio.
\item \textbf{No reinventar lo que ya funciona.} Los símbolos matemáticos,
 las fórmulas químicas y las partituras se escriben como siempre.
\item \textbf{Cada signo significa una sola cosa.} El tamaño, el color, el giro o
 la posición no cambian el significado. Un signo cambia de significado solo
 con una revisión escrita en el anexo de cambios.
\item \textbf{Las relaciones importantes no se mezclan.} Causa, sucesión,
 deducción, transformación y correlación son signos distintos. Negar, oponer y
 contrastar también (sección~\ref{sec:distinciones}).
\item \textbf{Lo que dice otro no es lo que pienso yo.} La fuente, mi propia idea,
 una hipótesis y una duda tienen marcas distintas.
\item \textbf{Crece con el uso.} Los signos nuevos salen de lo que escribo
 muchas veces, no de lo que suena bonito en teoría (sección~\ref{sec:crecer}).
\end{enumerate}

\subsection{Qué no es}
No es un idioma para hablar: no tiene pronunciación, y los nombres de los signos
son solo etiquetas para este manual. Tampoco es un lenguaje formal para
computadoras, ni un cifrado para que otros no entiendan. Los códigos
\texttt{K01}--\texttt{K35} sirven para buscar signos en el catálogo; no se
escriben en el cuaderno.
